% Bagian: computability
% Bab: computability-theory
% Subbagian: total

\documentclass[../../../include/open-logic-section]{subfiles}

\begin{document}

\olfileid[id]{cmp}{thy}{tot}
\olsection{Totalitas Tidak Dapat Diputuskan}

Mari kita tinjau satu contoh lagi tentang penggunaan teorema $s$-$m$-$n$
untuk menunjukkan bahwa sesuatu tidak dapat dihitung. Misalkan $\fn{Tot}$
himpunan indeks fungsi total yang dapat dihitung, yakni
\[
\fn{Tot} = \Setabs{x}{\text{bagi setiap $y$, $\cfind{x}(y)\fdefined$}}.
\]

\begin{prop}
\ollabel{prop:total}
$\fn{Tot}$ tidak dapat dihitung.
\end{prop}

\begin{proof}
Untuk melihat bahwa $\fn{Tot}$ tidak dapat dihitung, cukup ditunjukkan bahwa
$K$ tereduksi kepadanya. Definisikan $h(x,y)$ dengan
\[
h(x,y) \simeq
\begin{cases}
0 & \text{jika $x \in K$} \\
\fundefined & \text{selain itu.}
\end{cases}
\]
Perhatikan bahwa $h(x,y)$ sama sekali tidak bergantung pada~$y$. Tidak sulit
untuk melihat bahwa $h$ dapat dihitung parsial: pada masukan $x,y$, kita
menghitung~$h$ dengan mula-mula menyimulasikan fungsi~$\cfind{x}$ pada
masukan~$x$; jika komputasi ini berhenti, $h(x,y)$ mengeluarkan~$0$ lalu
berhenti. Jadi, $h(x,y)$ tidak lain adalah
$\Zero(\umin{s}{T(x,x,s)})$, dengan $\Zero$ fungsi nol konstan.

Menurut teorema $s$-$m$-$n$, terdapat fungsi rekursif primitif~$k(x)$
sedemikian sehingga bagi setiap $x$ dan~$y$,
\[
\cfind{k(x)}(y) \simeq
\begin{cases}
0 & \text{jika $x \in K$} \\
\fundefined & \text{selain itu.}
\end{cases}
\]
Jadi, $\cfind{k(x)}$ total jika $x \in K$, dan tidak terdefinisi pada setiap
masukan jika sebaliknya. Dengan demikian, $k$ merupakan reduksi dari $K$
ke~$\fn{Tot}$.
\end{proof}

\begin{digress}
Ternyata $\fn{Tot}$ bahkan tidak terenumerasi secara komputabel---kompleksitasnya
berada lebih tinggi dalam ``hierarki aritmetis.'' Namun, kita tidak akan
membahas penguatan ini di sini.
\end{digress}

\end{document}
